\subsection{Vektoren als Verschiebungspfeile}
\label{subsec:vector_algebra_vectors_as_displacements}

\begin{DefinitionBox}{Vektor}{vector_algebra_vector}
  Ein \emph{Vektor} \EnglishTerm{vector} ist durch seine Länge und seine
  Richtung bestimmt. Geometrisch stellen wir ihn als
  \emph{Verschiebungspfeil} \EnglishTerm{displacement arrow} dar und schreiben
  beispielsweise $\vec a$.
  \[
    \vec a
    \quad\widehat{=}\quad
    \text{Länge und Richtung}
  \]
  Ein Vektor besitzt keine feste Position im Raum: Paralleles Verschieben
  ändert den Vektor nicht.
\end{DefinitionBox}

\begin{center}
  \CachedTikzFigure{content/resources/chapters/vector_algebra/tikz/vector_translation.tex}
\end{center}

\begin{KeyPointBox}{Beziehungen zwischen Vektoren}{vector_algebra_vector_relations}
  \noindent
  \begin{minipage}[c]{0.50\linewidth}
    \begin{align*}
      \vec a=\vec b
      &\quad\Longleftrightarrow\quad
      \text{gleiche Länge und Richtung},\\[0.6em]
      \vec a\parallel\vec b
      &\quad\Longleftrightarrow\quad
      \text{parallel oder antiparallel},\\[0.6em]
      \vec b=-\vec a
      &\quad\Longleftrightarrow\quad
      \text{Gegenvektoren}.
    \end{align*}
  \end{minipage}\hfill
  \begin{minipage}[c]{0.46\linewidth}
    \centering
    \CachedTikzFigure{content/resources/chapters/vector_algebra/tikz/vector_relations.tex}
  \end{minipage}
\end{KeyPointBox}

Der \emph{Gegenvektor} \EnglishTerm{opposite vector} zu $\vec a$ ist $-\vec a$.
